\documentclass[12pt,a4paper]{article}

%% ============================================================
%%  Packages
%% ============================================================
\usepackage[margin=2.5cm]{geometry}
\usepackage{graphicx}
\usepackage[T1]{fontenc}
\usepackage[utf8]{inputenc}
\usepackage{booktabs}
\usepackage{longtable}
\usepackage{array}
\usepackage{makecell}
\usepackage{multirow}
\usepackage{amsmath}
\usepackage{amssymb}
\usepackage{enumitem}
\usepackage{parskip}
\graphicspath{{./}{figures/}}

%% ---- AVA brand style (colours, headings, boxes, listings) ------
\usepackage{ava_style}
\renewcommand{\AVAManualTitle}{System Manual}

%% ---- Header right-hand text (manual-specific) ------------------
\fancyhead[R]{\small\textcolor{avaInkBlue}{\textsc{System Manual}}}

%% ---- PDF metadata -----------------------------------------------
\hypersetup{
  pdftitle   = {AGN Variability Archive --- System Manual},
  pdfauthor  = {J.V. Hernandez Santisteban},
}

\begin{document}
%%============================================================

\avatitlepage
  {System Manual}
  {1.0}
  {16 November 2026}
  {Juan V.\ Hern\'{a}ndez Santisteban}
  {jvhs1@st-andrews.ac.uk}
  {School of Physics \& Astronomy, University of St Andrews,\\
   North Haugh, St Andrews, KY16~9SS, UK}

\tableofcontents
\clearpage

%%============================================================

\section{Overview}

\begin{notebox}
\textbf{Document scope.}  This manual covers system-level topics only:
architecture, database, configuration, orchestration, and installation.
For the photometric reduction chain (SExtractor, calibration, aperture
tables, inter-calibration) see the \emph{AVA Photometry Manual}
(\texttt{docs/photometry\_manual.tex}).  The spectroscopy reduction
manual is forthcoming pending completion of the improved banzai-floyds
pipeline.
\end{notebox}

\label{sec:overview}
%%============================================================

The AGN Variability Archive (AVA) photometric pipeline processes multi-band imaging light
curves from the Las Cumbres Observatory global telescope network for approximately 70 active
galactic nuclei.  The overarching scientific goal is reverberation mapping: measuring the
time lag between AGN accretion-disc continuum variability and the response in the broad
emission-line region, which yields black-hole mass estimates.

\subsection{Instruments and telescopes}

Photometry is obtained with two camera families:
\begin{itemize}
  \item \textbf{Sinistro} (fa-series detectors, 1\,m telescopes) at sites coj, ogg, lsc,
        cpt, elp, tfn.
  \item \textbf{Spectral} (fs-series detectors, 2\,m Faulkes telescopes) at ogg and coj.
\end{itemize}
Standard LCO broadband filters are used: $u'g'r'i'z'$ (SDSS system), plus $B$, $V$, $R$,
$I$ (Johnson--Cousins) for some targets.  FLOYDS spectrographs on the 2\,m telescopes
provide simultaneous spectroscopic monitoring handled by a separate pipeline
(see the FLOYDS RM report).  Some targets also have Liverpool Telescope (LT/IO:O)
photometry ingested through a parallel processing branch.

\subsection{Pipeline outputs}

For each target on each observing night the pipeline produces:
\begin{enumerate}
  \item Calibrated flux light curves (mJy, all filters), appended to \texttt{lco\_latest.pkl}.
  \item Publication-quality and diagnostic plots (\textsc{pdf} + \textsc{png}) copied to the
        AVA website.
  \item An inter-calibrated merged light curve correcting telescope-to-telescope flux offsets.
  \item Standard-star catalogues (\texttt{lco\_latest\_stan.pkl}) for internal photometric
        quality monitoring.
\end{enumerate}

\subsection{Data flow}

\begin{figure}[h!]
\centering
\begin{minipage}{0.85\linewidth}
\begin{lstlisting}[style=shell]
target_list.txt
      |
      v
ava_update.py  <-- flags: -o -d -e -c -p -m -s
  (one shell script per target, run with nohup)
      |
      +-> ava_initialise.py    set up directories, fetch APASS catalogue
      +-> download_data.py     pull BANZAI FITS from LCO archive API
      +-> ava_sextractor.py    SExtractor aperture photometry  ->  _cat.fits
      +-> ava_calibrate.py     match APASS, fit zeropoint  ->  lco_latest.pkl
      +-> ava_webplot.py       convert to mJy, plot all filters + Swift
      `-> ava_merge.py         CALI inter-calibration across telescopes
\end{lstlisting}
\end{minipage}
\caption{End-to-end data flow.  Each stage is an independent Python script
  invoked by a generated shell script; state passes via pickle files and FITS
  catalogues on disk.}
\label{fig:flow}
\end{figure}

\subsection{Technology stack}
Python~3, \texttt{numpy}/\texttt{scipy}, \texttt{astropy}, \texttt{pandas},
\texttt{matplotlib}, \texttt{lmfit}, SExtractor (external binary),
\texttt{pycali\_jvhs} (custom MCMC inter-calibration wrapper).

%%============================================================

\section{Known Issues, Limitations \& Migration Notes}
\label{sec:issues}
%%============================================================

This section catalogues the technical debt that must be resolved before the pipeline can
be safely migrated to a new machine or handed to a new user.

\begin{enumerate}

\item \textbf{Hardcoded absolute paths.}  Every script contains literal paths to
      \texttt{/data/jvhs1/agn/} or \texttt{/data3/jvhs1/agn/}.  These are the single
      largest portability blocker.  AVA~v2 must replace all of them with paths read from
      \texttt{config/ava\_config.yaml}.

\item \textbf{APASS via web scrape.}  \texttt{ava\_initialise.py} fetches the APASS
      catalogue using the \texttt{mechanize} browser library.  Replacement:
      \texttt{astroquery.vizier.query\_region()} on APASS\,DR9 or Pan-STARRS\,PS1.

\item \textbf{Pickle-based state.}  \texttt{lco\_latest.pkl} with nested arrays is not
      portable across versions.  Replace with HDF5 / Parquet or the \texttt{lcodb}
      SQLite database.

\item \textbf{SExtractor as external binary.}  \texttt{os.system(`sex ...')} with no
      version checking.  Replace with \texttt{sep} (Python SExtractor) or
      \texttt{photutils}.

\item \textbf{Aperture index by magic number.}  Bare integer literals with inline MJD
      comparisons.  Must become a named config parameter with documented physical meaning.

\item \textbf{Credential in source code.}  \texttt{ava\_update.py} contains a hardcoded
      LCO archive username and password.  Replace with the token-based auth used by
      \texttt{lcodb}.

\item \textbf{SCP for delivery.}  Hardcoded hostname and path for plot delivery.
      Decouple pipeline from delivery mechanism.

\item \textbf{pycali\_jvhs undocumented dependency.}  Private package not on PyPI.
      Must be packaged or replaced with PyROA / CREAM.

\item \textbf{No test suite.}  No automated tests for any stage.  AVA~v2 should add
      unit tests for calibration arithmetic and integration tests on reference frames.

\item \textbf{Parallel execution via shell.}  Per-target parallelism via
      \texttt{nohup}~scripts provides no job management, logging, or failure recovery.
      AVA~v2 should use Python \texttt{multiprocessing} or a lightweight task queue.

\end{enumerate}

%%============================================================

%%============================================================
\section{Installation Guide}
\label{sec:install}
%%============================================================

This chapter documents how to install the AVA pipeline on a new machine from
scratch.  Follow the steps in order.  Where a command differs between macOS
and Linux, both variants are shown.

\subsection{Prerequisites}

The pipeline requires the following to be present on the machine before
installation begins:

\begin{itemize}
  \item \textbf{Conda} (Miniconda or Anaconda) --- used to create an isolated
        Python environment.  Install from
        \url{https://docs.conda.io/en/latest/miniconda.html}.
  \item \textbf{Git} --- to clone the repositories.
  \item \textbf{SExtractor} --- the external source-extraction binary.
        On macOS: \texttt{brew install sextractor}.
        On Linux: \texttt{sudo apt install sextractor} or compile from source.
  \item \textbf{SSH access} to the LCO data storage machine (if working
        remotely).
\end{itemize}

\subsection{Step 1 --- Clone the repositories}

\begin{lstlisting}[style=shell]
# Choose a base directory (adjust to taste)
BASE=~/astro

mkdir -p $BASE && cd $BASE

# AVA main repository
git clone https://github.com/jvhs1/AVA.git

# banzai-floyds (spectroscopy pipeline)
# Use the rm-improvements branch for the RM-optimised version
git clone https://github.com/LCOGT/banzai-floyds.git floyds/banzai-floyds
cd $BASE/floyds/banzai-floyds
git checkout rm-improvements
cd $BASE
\end{lstlisting}

\subsection{Step 2 --- Create the conda environment}

\begin{lstlisting}[style=shell]
cd $BASE/AVA
conda env create -f environment.yml
conda activate ava
\end{lstlisting}

This creates an environment named \texttt{ava} with Python~3.12 and all
required packages.  To verify:
\begin{lstlisting}[style=shell]
python -c "import ava; print(ava.__version__)"
\end{lstlisting}

\subsection{Step 3 --- Install AVA and its sibling packages}

\begin{lstlisting}[style=shell]
# Install AVA itself in editable mode
pip install -e "$BASE/AVA"

# Install banzai-floyds in editable mode (spectroscopy only)
pip install -e "$BASE/floyds/banzai-floyds"

# pycali_jvhs (inter-calibration; install from local copy)
# pip install -e /path/to/pycali_jvhs
# Ask J.V.H.S. for the source if not already present.
\end{lstlisting}

After this step the \texttt{lcodb} command is available (it is bundled inside
AVA):
\begin{lstlisting}[style=shell]
lcodb --help
\end{lstlisting}

\subsection{Step 4 --- Set the LCO API token}

Obtain your personal API token from
\url{https://observe.lco.global} $\to$ Profile $\to$ API~token.

\begin{lstlisting}[style=shell]
echo 'YOUR_TOKEN_HERE' > ~/.lco_token
chmod 600 ~/.lco_token
\end{lstlisting}

\begin{warnbox}
Never put the token directly in a config file or commit it to git.
\end{warnbox}

\subsection{Step 5 --- Create and configure your local config file}

\begin{lstlisting}[style=shell]
cd $BASE/AVA
cp config/ava_config.yaml config/ava_config.local.yaml
\end{lstlisting}

Open \texttt{config/ava\_config.local.yaml} in an editor and update the
\texttt{paths} section to match the new machine.  The keys that \emph{must}
be set are listed in Table~\ref{tab:paths}.

\begin{table}[h!]
\centering
\small
\begin{tabular}{>{\ttfamily}p{5.2cm} p{4.5cm} p{4.0cm}}
\toprule
Key & Description & Production value \\
\midrule
paths.photometry\_root   & Root of photometry data tree & \texttt{/data3/jvhs1/agn} \\
paths.photometry\_raw    & Raw frames subfolder         & \texttt{raw\_data} \\
paths.spectra\_root      & Root of spectroscopy data tree & \texttt{/data2/jvhs1/agn} \\
paths.spectra\_raw       & Raw FLOYDS frames subfolder  & \texttt{raw\_data} \\
paths.plots\_dir         & Shared plot output directory & \texttt{/data3/jvhs1/agn/ava\_plots} \\
paths.swift\_db          & Swift/UVOT merged CSV        & \texttt{/data3/jvhs1/agn/ava\_swift\_updated.csv} \\
paths.template\_dir      & SExtractor config templates  & \texttt{/data3/jvhs1/agn/template\_pipeline} \\
paths.floyds\_pipeline   & banzai-floyds working copy   & \texttt{\$BASE/floyds/banzai-floyds} \\
paths.database           & AVA SQLite master DB         & \texttt{data/ava.db} \\
paths.ava\_csv           & AVA master CSV               & \texttt{data/ava/ava.csv} \\
\bottomrule
\end{tabular}
\caption{Path keys that must be set for each machine.  Relative paths are
  resolved from \texttt{ava\_root} (auto-detected as the repository root if
  left blank).}
\label{tab:paths}
\end{table}

Point AVA at the local config:
\begin{lstlisting}[style=shell]
# Option A: environment variable (add to ~/.bashrc or ~/.zshrc)
export AVA_CONFIG="$BASE/AVA/config/ava_config.local.yaml"

# Option B: symlink the local copy over the default
ln -sf ava_config.local.yaml config/ava_config.yaml
\end{lstlisting}

Verify that the config loads and all paths resolve:
\begin{lstlisting}[style=python]
from ava.config import cfg
print(cfg)
# AVAConfig(
#   config  = /home/.../config/ava_config.local.yaml
#   root    = /home/.../AVA
#   phot    = /data3/jvhs1/agn
#   spec    = /data2/jvhs1/agn
#   db      = /home/.../AVA/data/ava.db
# )
\end{lstlisting}

\subsection{Step 6 --- Initialise the LCO archive databases}

This step builds the \texttt{lcodb} archive databases by crawling the LCO
archive API.  It is a one-time operation that takes between a few minutes
(spectra) and several hours (photometry headers) depending on the data volume.

\begin{lstlisting}[style=shell]
# Crawl all KEY proposals (metadata only; no FITS files downloaded)
lcodb sync

# Download FITS headers for all spectra (~10^4 frames, minutes)
lcodb headers --db spectra

# Compute airmass for frames without it in the header
lcodb airmass

# Build science->standard and spectrum->arc/flat matching tables
lcodb match

# Link LCO target names to AVA sources
lcodb link

# Check status
lcodb status
\end{lstlisting}

After this step, the archive databases are ready at the paths set in
\texttt{config/lcodb\_config.yaml} (default: \texttt{data/lcodb/}).

\subsection{Step 7 --- Migrate the existing AVA database}

If migrating from an existing machine, copy the master SQLite database and
run the migration script to import it into the new schema:

\begin{lstlisting}[style=shell]
# Copy from the old machine (adjust path)
scp jvhs1@oldmachine:/path/to/ava.sqlite data/ava/ava.sqlite.old

# Migrate into the new AVA database
python -m ava.database.migrate \
    --source data/ava/ava.sqlite.old \
    --dest data/ava.db
\end{lstlisting}

\subsection{Step 8 --- Schedule the daily update}

The daily update fetches new frames from the LCO archive and updates the
databases.  Add the following to \texttt{crontab -e}:

\begin{lstlisting}[style=shell]
# Run every day at 09:30 UTC
30 9 * * * /path/to/AVA/scripts/lcodb_daily.sh
\end{lstlisting}

Or on macOS using \texttt{launchd}, create a plist pointing to the same script.

\subsection{Step 9 --- Verify the installation}

\begin{lstlisting}[style=python]
from ava.archive import get_archive
a = get_archive()
print(a)              # Archive(spectra=ok, photometry=ok, ava=ok)
a.overview()          # Should show the full AGN table
a.summary('Fairall 9')
\end{lstlisting}

Run the test suite to confirm all units pass:
\begin{lstlisting}[style=shell]
pytest tests/ -v
\end{lstlisting}

\subsection{Data directory structure on the production machine}

For reference, the expected layout on the production machine
(\texttt{/data2}, \texttt{/data3}) is:

\begin{lstlisting}[style=shell]
# Photometry  (--  /data3/jvhs1/agn/ --)
/data3/jvhs1/agn/
  raw_data/
    <target_slug>/
      YYYYMMDD/
        *e91.fits.fz        # BANZAI-reduced frames
  <target_slug>/
    apass_<target>.csv      # APASS reference catalogue
    apass_pytics.csv        # Reformatted for pytics
    lco_latest.pkl          # Calibrated photometry (all epochs)
    lco_latest_stan.pkl     # Calibrated standard stars
    lco_latest_*.pdf/png    # Light-curve plots
    sextractor_ava/
      *_cat.fits            # SExtractor catalogues
    ava_merging/            # CALI workspace
  ava_plots/                # Plots copied here before web delivery
  ava_swift_updated.csv     # Swift/UVOT merged database
  template_pipeline/        # SExtractor config templates

# Spectroscopy  (--  /data2/jvhs1/agn/ --)
/data2/jvhs1/agn/
  raw_data/
    <target_slug>/
      YYYYMMDD/
        *e91.fits.fz        # BANZAI FLOYDS frames
  <target_slug>/            # banzai-floyds reduction products
\end{lstlisting}


%%============================================================

\section{Spectroscopic Archive: Standards \& Coverage}
\label{sec:spec-archive}
%%============================================================

\textit{Last updated: 2026-11-15. Source:
\texttt{floyds/lco\_archive\_db/db/spectra.sqlite} (252\,MB).}

\subsection{Overall statistics}

\begin{table}[h!]
\centering
\small
\begin{tabular}{lr}
\toprule
Metric & Value \\
\midrule
Science spectra (raw, \texttt{rlevel=0}) & 9{,}524 \\
Standard spectra (raw) & 20{,}712 \\
Science spectra with $\geq 1$ matched standard & \textbf{9{,}524 (100\,\%)} \\
Science spectra with no standard & \textbf{0} \\
Total \texttt{std\_match} rows & 63{,}615 \\
BANZAI 1d products available & 3{,}110 / 21{,}477 (14.5\,\%) \\
Total FLOYDS exposure time & 1{,}926\,h across 8{,}860 frames, 5{,}139 epochs \\
Targets with any FLOYDS data & 57 of 71 \\
\bottomrule
\end{tabular}
\caption{FLOYDS spectroscopic archive summary (2026-11-15).}
\end{table}

\begin{notebox}
\textbf{Full standard coverage.}  Every one of the 9{,}524 science spectra has
at least one matched flux standard in \texttt{spectra.sqlite}.  No coverage
gaps exist in the current archive.
\end{notebox}

\subsection{Spectrograph and slit usage}

\begin{table}[h!]
\centering
\small
\begin{tabular}{llrll}
\toprule
Instrument & Site & Science spectra & First obs & Last obs \\
\midrule
en06 (FLOYDS-N) & ogg (Haleakala)     & 5{,}656 & 2014-08-05 & 2026-11-14 \\
en12 (FLOYDS-S) & coj (Siding Spring) & 2{,}570 & 2019-04-14 & 2026-11-11 \\
en05 (FLOYDS-S, retired) & coj        &   634   & 2014-05-19 & 2019-02-03 \\
\bottomrule
\end{tabular}
\caption{Science spectra by spectrograph.}
\end{table}

Slit-width distribution: 5{,}854 at 6\,arcsec; 1{,}380 at 2\,arcsec;
126 at 1.6\,arcsec; 10 at 1.2\,arcsec; 1{,}490 with no slit keyword in the
header.

\subsection{Standard selection}

The \texttt{spectrum\_standards} view in \texttt{spectra.sqlite} assigns one
standard per science spectrum using this priority order:

\begin{enumerate}
  \item \textbf{\texttt{same\_slit}} --- nearest standard taken with the same
        slit width (2{,}849 spectra, mean $\langle\Delta t\rangle = 176$\,h).
  \item \textbf{\texttt{nearest\_2arcsec}} --- nearest 2\,arcsec standard when
        no same-slit option exists (4{,}280 spectra, $\langle\Delta t\rangle =
        48$\,h).
  \item \textbf{\texttt{nearest\_any}} --- nearest standard of any slit
        (1{,}731 spectra, $\langle\Delta t\rangle = 30$\,h).
\end{enumerate}

Standards are \emph{never taken in the same observing block} as the science
(\texttt{same\_block = 0} throughout); FLOYDS standards are submitted as
separate queue requests.  The minimum $\Delta t$ is typically $\approx
12$\,minutes.

\subsection{Targets with large average standard $\Delta t$}
\label{sec:std-dt}

The targets below had sparse standard coverage during their campaigns and
may carry systematic flux calibration offsets.

\begin{table}[h!]
\centering
\small
\begin{tabular}{lrr}
\toprule
Target & Science spectra & $\langle|\Delta t|\rangle$ (hours) \\
\midrule
Mrk\,1220      &  86 & 552 \\
3C\,120        & 153 & 299 \\
Mrk\,110       & 217 & 232 \\
SDSS\,J2353    &  13 & 168 \\
PKS\,2004-447  &   7 &  88 \\
PG\,1302-102   &  99 &  71 \\
NGC\,3783      & 177 &  58 \\
Mrk\,1239      &  49 &  55 \\
Mrk\,766       &  50 &  50 \\
\bottomrule
\end{tabular}
\caption{Targets with $\langle|\Delta t|\rangle > 50$\,h to nearest standard.
  These are predominantly old campaigns (KEY2014A-002, KEY2018B-001) where
  dedicated FLOYDS standards were taken less frequently.}
\end{table}

\subsection{Targets with no FLOYDS data (14)}

SDSS-RM201, SDSS-RM562, PG\,1310-108, 2MASS\,J21368-6438,
2MASX\,J21075970-3754087, PG\,2304+043, RX\,J22198-0935, NGC\,4395,
PKS\,0558-504, NGC\,2617, Zw\,535-12, SDSS\,J1430+2303, Mrk\,704,
IGR\,J17062-6143.

\subsection{Database update log}

\begin{table}[h!]
\centering
\small
\begin{tabular}{lll}
\toprule
File & Location & Updated \\
\midrule
\texttt{ava.db}                 & \texttt{AVA/data/} & 2026-11-15 \\
\texttt{ava.csv}                & \texttt{AVA/data/ava/} & 2026-11-15 \\
\texttt{ava\_link\_overrides.csv} & \texttt{AVA/data/ava/} & 2026-11-15 \\
\texttt{spectra.sqlite}         & \texttt{floyds/lco\_archive\_db/db/} & 2026-11-15 \\
\texttt{photometry.sqlite}      & \texttt{floyds/lco\_archive\_db/db/} & live (daily) \\
\bottomrule
\end{tabular}
\caption{AVA data files and their last update dates.  \texttt{spectra.sqlite}
  and \texttt{photometry.sqlite} remain in the \texttt{floyds/} sandbox and
  are accessed read-only via \texttt{get\_archive()}.  A snapshot of
  \texttt{spectra.sqlite} taken before the 2026-11-15 calibration update is
  kept as \texttt{spectra\_before\_calib\_20261115.sqlite}.}
\end{table}



\appendix
\section{Running the Pipeline}
\label{app:run}
%%============================================================

\textbf{Full run for a single target on the latest data:}
\begin{lstlisting}[style=shell]
python ava_update.py -d -e -c -p -list target_list.txt -project latest
\end{lstlisting}

\textbf{Monitor progress:}
\begin{lstlisting}[style=shell]
ps aux | grep 'python download_'
ps aux | grep 'python ava_sextractor'
ps aux | grep 'python ava_calibrate'
tail -100 lco_Fairall_9_logall.out
\end{lstlisting}

\textbf{Force full recalibration (e.g.\ after aperture change):}
\begin{lstlisting}[style=shell]
python ava_calibrate.py -obj Fairall_9 -s -j   # -j applies 2026 aperture fix
\end{lstlisting}

\textbf{Inter-calibration:}
\begin{lstlisting}[style=shell]
python ava_update.py -m -s -list target_list.txt
\end{lstlisting}

%%============================================================
\section{Glossary}
\label{app:glossary}
%%============================================================

\begin{description}
  \item[APASS] AAVSO Photometric All-Sky Survey.  Provides $BVgri$ magnitudes for
        stars used as photometric standards.
  \item[BANZAI] LCO automated image-reduction pipeline.  Delivers bias-subtracted,
        flat-fielded, astrometry-solved FITS frames at \texttt{rlevel=91}.
  \item[CALI] Cross-calibration algorithm for multi-telescope light curves.  Fits
        per-telescope flux scale factors using MCMC.
  \item[MJD] Modified Julian Date. MJD = JD $-$ 2400000.5.
  \item[ZP] Photometric zeropoint.  The magnitude of a source that would produce
        1 ADU/s in the instrument.
  \item[SExtractor] Source-Extractor; software for source detection and aperture
        photometry in astronomical images.
\end{description}

\vfill
\hrule
\vspace{4pt}
{\small
Created: 2026-11-15 \quad
Author: J.V.\ Hernandez Santisteban \quad
Pipeline version: v1.3 (legacy) / AVA~v2 (in development)
}

\end{document}